\Subhead{1.6.1}{Автоколебательные системы с двумя и более степенями свободы}
В активной системе с несколькими степенями свободы возможны несколько нормальных мод. Например, LC-контур генератора связан со вторым пассивным контуром: у пары две собственные частоты, но энергию обе моды получают через один транзистор. Насыщение одной модой меняет усиление другой --- возникает \emph{конкуренция мод}. Линейное разложение не превращает нелинейную систему в независимые генераторы.

При слабой нелинейности и отсутствии близкого комбинационного резонанса используем
\begin{equation}\label{eq:modecompetition}
 \dot A_1=A_1(\mu_1-b_{11}A_1^2-b_{12}A_2^2),\quad
 \dot A_2=A_2(\mu_2-b_{21}A_1^2-b_{22}A_2^2).
\end{equation}
$\mu_j$ --- малосигнальные инкременты; $b_{jj}>0$ --- самонасыщение, $b_{jk}$ --- перекрёстное. Последние определяются общим активным элементом и перекрытием мод. Для $I_j=A_j^2$ совместный режим удовлетворяет
\begin{equation}
 b_{11}I_1+b_{12}I_2=\mu_1,\quad
 b_{21}I_1+b_{22}I_2=\mu_2.
\end{equation}
Сначала нужны $I_1,I_2>0$, затем проверка устойчивости. При этих условиях совместный режим устойчив по интенсивностям при $D=b_{11}b_{22}-b_{12}b_{21}>0$; при $D<0$ он седловой. Одномодовый режим $I_1=\mu_1/b_{11}>0$ устойчив к появлению второй моды, если $\mu_2-b_{21}\mu_1/b_{11}<0$. Оба одномодовых режима могут быть устойчивы: тогда возможны бистабильность и выбор по предыстории.

При перестройке пассивного контура меняются частоты и доля каждой моды в активном контуре. Генерация может следовать одной ветви избегаемого пересечения, пока усиления хватает, затем перескочить на другую. Гистерезис возникает, если области устойчивости перекрываются; он не следует из одного лишь расщепления линейных частот. Такая последовательность соответствует обсуждению связанных контуров в конспекте Балакшия (PDF-с.~16, лист~14). Для лазера та же логика рассматривается в~\Qref{1.10.2}.
